\section{One-Loop Factorization}

Now we are ready to construct a factorization formula at one-loop order. In the infinite momentum frame or on the light-cone, the momentum fraction in
parton distributions and splitting functions is limited to $[0,1]$. However, in the present case, the splitting in
the quasi-distribution is not constrained to this region, it can be in $[-\infty,\infty]$. Thus the connection
of the two distributions is reflected through the following factorization theorem up to power corrections in the large $P^z$ limit
\begin{equation}
    \tilde q(x, \mu^2, P^z)  = \int^1_{0} \frac{dy}{y} Z\left(\frac{x}{y},\frac{\mu}{P^z}\right) q(y, \mu^2) +  {\cal O}\left(\Lambda^2/(P^z)^2,  M^2/(P^z)^2\right)\ ,
\end{equation}
where the integration range is determined
by the support of the quark distribution $q(y)$ on the light cone.

The $Z$ factor has a perturbative expansion in $\alpha_s$
\begin{equation}
             Z(\xi, P^z/\mu) = \delta (\xi-1) + \frac{\alpha_s}{2\pi} Z^{(1)}(\xi, P^z/\mu) + \dots
\end{equation}
Before we proceed, it is important to note the existence of a linear divergence coupled to the axial-gauge singularity, which is a double pole at $\xi=1$, in the one-loop corrections to the quasi quark distribution, as one can see from Eqs.~(\ref{unpolvertex}) and (\ref{unpolself}) above. As mentioned in the previous section, if one chooses a covariant gauge like the Feynman gauge, the diagrams in Fig.~\ref{onelooppdf} do not give axial singularities, but now one has extra diagrams with gauge link and the singularities come from these diagrams. Obviously the double pole structure is absent in dimensional regularization, but it is present in the cutoff regularization used here. Unlike in the case of light-cone parton distribution, where one encounters at most a single pole at $\xi=1$ and can appropriately regularize it by a plus-prescription, the plus-prescription is unable to fully regularize the divergence of the double pole. In general, it only reduces the double pole to a single pole, which still yields singular contribution upon integration over $\xi$. However, it turns out that in our case this singularity can be removed with a principal value prescription, which corresponds to a regularization of the Wilson line self-energy. Although the linear divergence can also be singled out (this can be done on the lattice by varying the lattice spacing with fixed $P^z$ and $x$) and subtracted~\cite{footnote1}, we propose to include this linearly divergent term within our factorization formula, in order to simplify the extraction of light-cone distribution from lattice data.


Now we are ready to write down the matching factor connecting the quasi quark distribution to the light-cone quark distribution. For $\xi>1$, one has
\begin{equation}
   Z^{(1)}(\xi)/C_F = \left(\frac{1+\xi^2}{1-\xi}\right)\ln \frac{\xi}{\xi-1} + 1 + \frac{1}{(1-\xi)^2}\frac{\mu}{P^z}  \ ,
\end{equation}
whereas for $0<\xi<1$
\begin{equation}
   Z^{(1)}(\xi)/C_F = \left(\frac{1+\xi^2}{1-\xi}\right)\ln \frac{(P^z)^2}{\mu^2}+\left(\frac{1+\xi^2}{1-\xi}\right)\ln\big[{4\xi(1-\xi)}\big]-\frac{2\xi}{1-\xi}+1+\frac{\mu}{(1-\xi)^2P^z} \ ,
\end{equation}
and for $\xi<0$
\begin{equation}
   Z^{(1)}(\xi)/C_F = \left(\frac{1+\xi^2}{1-\xi}\right)\ln \frac{\xi-1}{\xi}-1 +\frac{\mu}{(1-\xi)^2P^z} \ .
\end{equation}
Near $\xi=1$, one has an additional term coming from the self energy correction
\begin{equation}
    Z^{(1)}(\xi) = \delta Z^{(1)}(2\pi/\alpha_s) \delta(\xi-1)\ ,
\end{equation}
which can be extracted from Eqs.~(\ref{unpolself}) and (\ref{unpolselfIMF}) above, and provides a plus-regularization for the singularity at $\xi=1$. Thus the large logarithmic dependence on $P^z$ in $\tilde q(x,\mu^2,P^z)$ can be transformed into the
renormalization scale dependence through the above matching condition.
On the lattice, the matching must be recalculated up to a constant accuracy
using the standard approach~\cite{ren}.

So far, we have considered only the quark contribution. One can start with an antiquark to do the one-loop
calculation. In this case, one also has a contribution to $\tilde q(x, \mu^2, P^z)$ from $\bar q(x)$. However,
the antiquark distribution has the property
\begin{equation}
     \bar q(x) = - q(-x)\ ,
\end{equation}
which is related to quark distribution at negative $x$.
By including both quark and antiquark contribution, one obtains the following factorization with the integration extended from $-1$ to $+1$,
\begin{equation}\label{1loopfact}
    \tilde q(x, \mu^2, P^z)  = \int^1_{-1} \frac{dy}{|y|} Z\left(\frac{x}{y},\frac{\mu}{P^z}\right) q(y, \mu^2) +  {\cal O}\left(\Lambda^2/(P^z)^2,  M^2/(P^z)^2\right)\ ,
\end{equation}
where the negative $y$ region includes the antiquark contribution. The above is the complete one-loop factorization theorem, which replaces
Eq. (11) and the $Z$-factor in Ref. \cite{Ji:2013dva}.

We have constructed at one-loop level a factorization formula connecting the quasi parton distribution to the light-cone parton distribution. Of course it remains to be shown that there exists such a formula to all-loop orders. The factorization formula then allows one to extract the parton distribution $q(x,\mu^2)$ from calculating the quasi parton distribution on the lattice by measuring
the time-independent, non-local quark correlator $\overline{\psi}(z)\gamma^zL(z,0)\psi(0)$ with
$L(z,z_0) = \exp\left(-ig\int^{z}_{z_0} dz' A^z(z') \right)$ in a state with increasingly large $P^z$
(maximum $\sim 1/a$ with $a$ denoting lattice spacing).
